\chapter*{About these notes}
\addcontentsline{toc}{chapter}{About these notes}

These are lecture notes for IEOR 222, \emph{Financial Engineering}, at UC Berkeley: the
mathematics of pricing and hedging derivatives, from discrete-time models to stochastic
calculus, risk-neutral pricing and the numerical methods used in practice.

% TODO(Alberto): check the wording of this acknowledgement before publishing.
They grew out of the discussion sections I taught as graduate student instructor for the
course, which is taught by Thibaut Mastrolia, and they follow the structure of his
lectures. The text, the proofs, the demos and the wording of the exercises were written
for this repository; any mistakes are mine.

\paragraph{How to use them.} Each chapter has three parts. The notes, in this document,
develop the theory with complete proofs wherever the argument teaches something, and
with precise references where it does not. The exercises are collected in a separate
document; solutions are not published, so that the course is not spoiled for the next
class. The interactive demos, marked \textcolor{cobalt}{$\blacktriangleright$}, are Python
notebooks that open in Google Colab in one click: they let you move the parameters of
each result and see it happen.

\paragraph{Extras.} Part~V collects topics I would teach if the semester were longer:
arbitrage-free option surfaces built from market quotes, Hawkes processes and limit order
books, rough volatility, and signatures. They assume the first four parts.

\paragraph{Prerequisites.} Undergraduate probability (independence, the central limit
theorem, Gaussian vectors, characteristic functions), multivariable calculus, and enough
Python to read short scripts. A few results from measure theory are used without proof
and stated where they first appear: the monotone and dominated convergence theorems, the
Borel--Cantelli lemma, Dynkin's $\pi$--$\lambda$ theorem, and the Tonelli and Fubini
theorems for exchanging integrals. Beyond these, no measure theory is assumed: the
measure-theoretic ideas the course needs are introduced in \cref{ch:01-probability}.

\paragraph{Corrections.} Please report errors, typos and unclear passages as issues on
the \href{\repourl/issues}{GitHub repository}.

\section*{Notation}
\begin{center}
\begin{tabular}{@{}ll@{}}
\toprule
$(\Omega,\F,\Prob)$ & probability space; $\Prob$ is the real-world (historical) measure\\
$\Q$, $\EQ$ & a risk-neutral measure and its expectation\\
$(\F_t)_{t\ge0}$ & filtration: the information available at time $t$\\
$W$ & a Brownian motion\\
$\mathcal N(m,v)$, $N(\cdot)$ & normal law with mean $m$ and variance $v$; standard normal CDF\\
$\N$ & $\{0,1,2,\dots\}$\\
$x^+$, $x\wedge y$ & $\max(x,0)$; $\min(x,y)$\\
$\1_A$ & indicator of the event $A$\\
$[X,X]_t$ & quadratic variation of $X$ on $[0,t]$\\
$r$, $\sigma$, $K$, $T$ & interest rate, volatility, strike, maturity\\
\bottomrule
\end{tabular}
\end{center}
